% 本文件由 LaTeX 排版 Agent 根据有效 translation.json 自动生成。
% author=Lactantius; work_id=0705; title=论迫害者去世的方式

\chapter{论迫害者去世的方式}

% structure: p0001 source=h1 target=omitted reason=page-title-already-rendered-by-parent

致\Person{多纳图}\par

% structure: p0003 source=h2 target=section reason=relative-heading-rank; base=h2

\section{第一章}

主已垂听了您，我最亲爱的多纳图，终日在他面前倾流的恳求，也垂听了我们其余弟兄的恳求，他们因光荣的宣认而获得了不朽的荣冠，即他们信德的赏报。请看，所有的仇敌都被消灭，整个罗马帝国重获安宁，那曾受压迫的教会再次兴起，被恶人之手推倒的天主圣殿，以比从前更大的荣耀重建起来。因为天主兴起了君王，废除暴君们不虔敬而血腥的法令，为人类谋求福祉；如此，过去的乌云消散，和平与宁静使众人心喜。在狂烈的旋风与黑暗的风暴之后，如今天空变得平静，渴望已久的光明已照耀出来。如今，天主，那听祈祷者，以他的神圣助佑，将他仆倒在地、受折磨的仆人从地上扶起，终止了恶人的合谋，擦去了哀悼者脸上的眼泪。那些侮辱天主者，躺卧在低处；那些推倒圣殿者，以更可怕的毁灭倒下；正义之人的折磨者，在他们为天主所降的瘟疫和应受的酷刑中，倾出了他们罪恶的灵魂。因为天主延迟惩罚他们，为要借伟大而奇妙的榜样，教导后世：惟独他是天主，并以适当的报复，对骄傲、不敬和迫害者施行审判。\par

关于这些人的结局，我以为宜于发表一篇叙述，使所有远方的人，以及所有以后兴起的人，都能知道全能者如何彰显他的权能和至高伟大，将仇视他名的人连根拔除，彻底消灭。这一点将变得明显，当我叙述从教会最初建立以来，谁是教会的迫害者，以及天主的审判者以何等严厉的惩罚向他们施行报应。\par

% structure: p0006 source=h2 target=section reason=relative-heading-rank; base=h2

\section{第二章}

在提庇留皇帝末年，鲁贝里乌斯·格米努斯与富菲乌斯·格米努斯执政期间，并在四月朔日前的第十日，如我所见记载，耶稣基督被犹太人钉在十字架上。他在第三日复活后，召集了他的宗徒们——他们在他被拿住时因恐惧而逃散——并同他们居住四十天，开启他们的心，为他们阐释至今仍被晦暗包裹的圣经，任命并装备他们宣讲他的圣言和教训，并规范了关于新约制度的一切。这一切完成后，一片云彩和旋风将他裹住，从人的视线中提到天上去了。\par

\Emph{（四月朔日前的第十日）}\par

学者们在使\Person{拉克坦提乌斯}自洽方面遇到了严重困难，如果这错误不是出于抄写者而是他自己的话。在\WorkTitle{训导}第四卷中，\Person{巴吕兹}给出的文字如下：\par

\Quote{在提庇留皇帝末期，正如我们所读到的记载，我们的主耶稣基督在双子执政官之年，四月朔日后的第十日，被犹太人钉在十字架上。}\par

\Person{拉克坦提乌斯}在\Place{尼科美底亚}写作，可能凭记忆引用了所读到的内容，也许是在\Person{比拉多}本人的报告中。\Emph{四月朔日后的第十日}这一表述是模糊的：\Person{贾维斯}说：\Quote{我的印象是，它意思是\InnerQuote{四月朔日前第十日之后}；也就是说，三月二十三日之后。}\par

但在这里，根据\Person{瓦尔丘斯}（公元1715年）的精确版本，我们的作者说——\par

\Quote{从此他们拥有分封侯，直到\Person{黑落德}，他处于\Person{提庇留}皇帝治下；在该皇帝第十五年，即双子执政官之年，四月朔日前第七日，犹太人将基督钉在十字架上。}\par

但在这里，根据四十份手稿的权威，\Person{杜·弗雷斯诺}读到\Quote{四月朔日前第十日}，他努力使之与上文的\Quote{四月朔日后第十日}协调。\Person{贾维斯}坚持采用\Quote{前第七日}的读法，有超过五十份手稿支持，并判定为三月二十三日。\par

他在著名的段落中引用了\Person{奥斯定}同样的意见：\par

\Quote{然而，基督是在那个月受孕并被钉死的，这由逾越节的遵守以及教会极熟知的他的诞辰日所表明。因为他是在第九个月出生，即一月朔日前第八日，那么他确实是在第一个月受孕，大约在四月朔日前第八日，这正是他受难的时间。}\par

\Person{奥斯定}认为这是按神秘意义\Quote{用母亲的奶煮山羊羔}，残忍地使十字架与童贞玛利亚的母职相重合，她在天使向她请安的同周年纪念日，目睹了自己的儿子成为无辜的牺牲品。\par

那时，他的宗徒共有十一人，后来加上了\Person{玛弟亚}，替代了叛徒\Person{犹达斯}，之后又加上了\Person{保禄}。然后，他们分散到全地，传扬福音，正如他们的主和师傅所命令的；在二十五年间，直到\Person{尼禄}皇帝统治之初，他们致力于在各行省和城市建立教会的基础。当\Person{尼禄}统治时，宗徒\Person{伯多禄}来到罗马，藉着天主所赐给他的能力，行了一些奇迹，并使许多人归向真宗教，为主建立了一座忠实而稳固的殿宇。\Person{尼禄}听到这些事，并看到不仅在罗马，而且在每一个地方，都有大批群众每日离弃偶像崇拜，谴责他们旧有的方式，转而皈依新宗教，他，一个可憎而凶残的暴君，便起身要摧毁这天上的殿宇并消灭真信仰。他是首先迫害天主仆人的人；他钉死了\Person{伯多禄}，杀死了\Person{保禄}：他也没有逃脱惩罚；因为天主垂顾了他子民的苦难；因此，这暴君被剥夺了权柄，从帝国的顶峰坠落，忽然消失，就连那恶兽的葬身之地也无处可寻。这使一些想象力丰富的人以为，他被带到了一个遥远的地区，仍然活着；他们把有关这事的\WorkTitle{西卜林}诗句应用在他身上：\par

\Quote{那杀死自己母亲的逃亡者，要从地极而来；}\par

仿佛那第一个迫害者也将成为最后一个，从而成为\OriginalTerm{敌基督}{Antichrist}的先行者！但我们不应相信那些人，他们断言两位先知\Person{哈诺客}和\Person{厄里亚}已经被迁到某个遥远的地方，以便在主来审判时侍奉他，同时幻想\Person{尼禄}将来要作为魔鬼的先行者出现，来毁灭大地和推翻人类。\par

% structure: p0021 source=h2 target=section reason=relative-heading-rank; base=h2

\section{第三章}

\Person{尼禄}死后过了几年，又兴起了一个同样邪恶的暴君（\Person{多米提安}），尽管他的统治极其可憎，但他长久地压迫他的臣民，并安坐王位，直到最后他伸出不敬的手攻击主。他受恶鬼鼓动去迫害义人，随即被交在敌人手中，受了当得的惩罚。在他自己的宫殿中被杀，这报复还不够：就连他的名号也被抹去。因为他虽然建筑了许多宏伟的建筑，重建了\Place{卡皮托利}，并留下了其他显赫的功绩，但元老院如此迫害他的名号，以至于他的雕像或纪念他的铭文痕迹全无；并且以最庄严、最严厉的敕令，在他死后也给他打上永久的耻辱烙印。这样，暴君的命令被撤销后，教会不仅恢复到原先的状态，而且更加光辉灿烂，日益兴旺。在随后的时代，当许多贤德的君主执掌罗马帝国的舵盘时，教会没有受到敌人猛烈的攻击，她将双手伸向东方和西方，以至于地上再也没有任何最遥远的角落没有传到这神圣的宗教，也没有任何习俗野蛮的民族，在转向崇拜天主后，没有变得温良与驯服。\par

% structure: p0023 source=h2 target=section reason=relative-heading-rank; base=h2

\section{第四章}

然而，这长久的和平后来被中断了。\Person{德西乌斯}出现在世上，一条可咒的恶兽，来折磨教会——若非坏人，谁会迫害宗教呢？他似乎是蒙召登上至尊之位，一面同时向天主发怒，一面同时倒下；因为他率军出征\OriginalTerm{卡尔皮人}{Carpi}，当时他们已占据了\OriginalTerm{达契亚}{Dacia}和\OriginalTerm{莫西亚}{Moefia}，他突然被蛮族包围，与他的大部分军队一同被杀；他甚至得不到葬礼仪式的尊荣，而是被剥光衣服，赤身露体，躺在地上成为野兽和飞鸟的食物——这是天主之敌的应得下场。\par

% structure: p0025 source=h2 target=section reason=relative-heading-rank; base=h2

\section{第五章}

紧接着，\Person{瓦莱里安}也以同样疯狂的姿态，举起不敬的手来攻击天主，尽管他统治时间不长，却流了许多义人的血。但天主以一种新颖而特别的方式惩罚了他，好作为后世的教训：天主的仇敌总要领受他们罪孽的公正报应。他被波斯人俘虏后，不仅失去了他滥用过度的权力，也失去了他剥夺他人自由的自由，并在最卑贱的奴隶状态下度过余生；因为波斯王\Person{沙普尔}，就是俘虏他的那位，无论何时要上车或上马，都命令那罗马人弯下腰来，献上他的背；然后，他把脚踩在\Person{瓦莱里安}的肩上，带着嘲笑的微笑说：\Quote{这才是真实的，而不是罗马人在木板或墙壁上描绘的那种。}\Person{瓦莱里安}在胜利者的当受侮辱中生活了相当长的时间；以致罗马之名长久成为蛮族的嘲笑和讥讽的对象；这在严厉的惩罚上又加上了一条：虽然他有一个做皇帝的儿子，却没有人来报复他的被俘和极其卑贱的奴隶状态；而且他也从未被要求归还。后来，他在如此大的耻辱中结束了这可耻的一生后，被剥了皮，他的皮从肉体上剥离，用朱砂染红，放在蛮族的神庙中，好让如此显赫的胜利的记忆得以永存，并让这景象常常展示给我们的使者，作为对罗马人的告诫：当他们看到被俘皇帝的战利品放在波斯神庙中时，不应过分自信于自己的力量。\par

既然天主如此惩罚了亵渎者，那么此后还有人胆敢对那唯一、统管并维系万有的天主的威严做出或甚至图谋任何事，岂不奇怪吗？\par

% structure: p0028 source=h2 target=section reason=relative-heading-rank; base=h2

\section{第六章}

\Person{奥勒良}或许能记起那位被俘皇帝的命运，但他天性暴虐固执，既忘了自己的罪，也忘了它的惩罚，竟以残酷的作为激起了天主的义怒。然而，他并未获准完成他的谋划；正当他开始放纵愤怒时，便被杀了。他的血腥敕令尚未送达较远的行省，他自己已浑身鲜血倒在\Place{色雷斯}的\Place{凯诺弗鲁里翁}，被他亲近的朋友刺杀，这些人对他产生了毫无根据的猜疑。\par

如此众多的此类例子本应使后来的暴君望而却步；然而他们不仅不惊惧，反倒在得罪天主的恶行上更加胆大妄为。\par

% structure: p0031 source=h2 target=section reason=relative-heading-rank; base=h2

\section{第七章}

当\Person{戴克里先}，那个邪恶的创造者、苦难的策划者，正在毁灭一切时，他甚至连天主也不放过侮辱。此人，部分由于贪婪，部分由于怯懦的计谋，颠覆了罗马帝国。因为他挑选了三人与他分享统治权；这样一来，帝国被分成四份，军队倍增，四个君主中的每一个都竭力维持着比以往任何单独皇帝都要庞大得多的军事力量。纳税的人开始比领薪饷的人还要少；因此，农民的资财被沉重的赋税榨干，农田被废弃，耕地变成林地，普遍恐慌蔓延。此外，行省被分割成细小的部分，许多总督和大量下级官员压在每一片领土、几乎每一座城市上。还有许多不同等级的管家和总督的副手。很少有民事案件提交给他们；但每天都有定罪，频繁的没收；无数商品的税，不仅经常重复，而且是永久的，在征收时还有难以忍受的暴行。\par

凡是为供养士兵而征收的，或许还可以忍受；但\Person{戴克里先}因其无法满足的贪婪，从不允许他国库里的钱减少；他不断搜集额外的补助和馈赠，以便他最初的库存保持不动和不可侵犯。他还通过种种勒索使万物变得极其昂贵后，试图通过法令来限制价格。于是为了最微不足道的小事流了许多血；人们害怕出售任何东西，匮乏变得比以往更加严重和悲惨，直到最后，这道法令在毁灭了无数人之后，出于纯粹的需要而被废除。此外，还有一种无尽的建筑热情，为了这个原因，从各省不断地勒索来支付工人和工匠的工资，并提供车辆及他计划中的工程所需的一切。这里建公堂，那里建竞技场，这里造铸币厂，那里造兵工厂；一处为他皇后建府邸，另一处为他的女儿建府邸。不久，城市大部分被遗弃，所有人都带着妻儿搬迁，就像从被敌人攻占的城镇搬走一样；当这些建筑建完，毁掉了整个省份之后，他说：\Quote{它们不对，要按另一方案重做。}于是它们又被拆除或改建，或许还会再次被拆毁。通过如此愚蠢的行为，他不断试图使\Place{尼科美底亚}在宏伟上与\Place{罗马}城媲美。\par

我略过不提有多少人因财产或财富而丧命；因为这类恶行极为频繁，并且因频繁而显得几乎合法。但这是他的特点：每当他看到一块田地格外精耕细作，或一栋房子异常雅致，立即就对该物主罗织罪名并准备处以死刑；所以，似乎\Person{戴克里先}若不流血，就不能犯下掠夺之罪。\par

% structure: p0035 source=h2 target=section reason=relative-heading-rank; base=h2

\section{第八章}

他的帝国兄弟\Person{马克西米安}（被称为\OriginalTerm{赫尔库利乌斯}{Herculius}）的性格如何？与\Person{戴克里先}并非不同；事实上，要使他们的友谊如此亲密和忠诚，他们必须有相同的倾向和目的，一致的意志和判断上的同心。他们只有这一点不同：\Person{戴克里先}更贪婪而更不果断；而\Person{马克西米安}贪婪较少，却更为大胆，不是倾向于善，而是倾向于恶。因为当他拥有\Place{意大利}——帝国的主要所在地——以及其他非常富饶的行省如\Place{阿非利加}和\Place{西班牙}近在咫尺时，他很少费心去保全那些他轻易就能积聚的财宝。每当他需要更多时，最富有的元老立即受到指使的证人的指控，说他们有篡位之罪；因此，元老院的主要人物天天被清除。于是，嗜血的国库充满了不义之财。\par

此外还有这个恶贯满盈者的不节制，不仅奸淫男性（这是可憎可恶的），而且奸污国家重要人物的女儿；凡他所到之处，少女突然从父母面前被夺走。在这些暴行中他追求至高快乐，而尽量满足他的淫欲和邪念在他看来便是他统治的幸福。\par

我略过\Person{君士坦提乌斯}不提，他是一位与众不同的君主，本应独揽帝国大权。\par

% structure: p0039 source=h2 target=section reason=relative-heading-rank; base=h2

\section{第九章}

但另一位马克西米安（\Person{伽列里乌斯}），被\Person{戴克里先}选为女婿，不仅比我们当代所经验的两位君主更坏，而且比往昔所有坏君主都更坏。在这头野兽内里，住着一种天然的野蛮和罗马血统所无的凶残；这不足为奇，因为他的母亲生在多瑙河对岸，是\Place{卡尔皮人}的入侵迫使她渡河逃到\Place{新达契亚}避难。\Person{伽列里乌斯}的相貌与他的作风相符：身材高大、肌肉丰满，肿胀到可怕的大块头；他的言语、姿势和眼神，使所有接近他的人都感到恐惧。他的岳父也极其惧怕他。原因如下：\Person{波斯王纳尔塞斯}效法他的祖父\Person{沙普尔}的榜样，集结大军，企图成为\Place{罗马帝国}东部各省的主人。\Person{戴克里先}易于在每次动乱中消沉胆怯，又担心遭遇像\Person{瓦勒良}那样的命运，不愿亲自迎战\Person{纳尔塞斯}；但他派\Person{伽列里乌斯}取道\Place{亚美尼亚}，而自己则停留在东部各省，焦虑地等待结果。野蛮人的习俗是携带一切属于他们的东西上战场。\Person{伽列里乌斯}设下伏兵，轻易击溃了因众多随从和辎重而陷入困境的敌人。他赶走了\Person{纳尔塞斯}，携带大量战利品归来，这大大助长了他的骄傲和\Person{戴克里先}的恐惧。因为此役之后，他傲慢到极点，竟拒绝凯撒的头衔；当他听到信函中以这个头衔称呼他时，便厉色凶声喊道：\Quote{我还要当凯撒多久？}于是他开始行为乖张，以至于好像他是另一个\Person{罗慕路斯}，希望被当作\Person{玛尔斯}的后裔并被称为如此；为了显得是神的后代，他情愿让自己的母亲\Person{罗慕拉}蒙受奸妇的恶名。但是，为了不混淆事件的年代顺序，我暂缓叙述他的事迹；因为后来当\Person{伽列里乌斯}获得皇帝称号，他的岳父被剥夺了帝位后，他变得完全疯狂，傲慢无度。\par

当\Person{狄奥克莱斯}——这是\Person{戴克里先}在取得君权之前的名字——以这样的行为并有这样的同伙致力于颠覆共和时，他的罪行理应承受一切灾祸；然而，只要他避免用义人的血玷污自己的双手，他统治得最为繁荣；至于他有何理由迫害义人，我现在就来解释。\par

% structure: p0042 source=h2 target=section reason=relative-heading-rank; base=h2

\section{第十章}

\Person{戴克里先}生性胆怯，喜好探求未来；在东方居住期间，他开始宰杀祭牲，企图从它们的内脏预知未来。当他献祭时，一些在场的基督徒侍从在额上画了不朽的标记。因此魔鬼被驱散，圣礼中断。占卜者战栗不已，无法在祭牲内脏上查得惯常的征兆。他们屡次重复献祭，仿佛先前的献祭不吉利；但一次次被宰杀的祭牲不提供任何占卜标记。最后，占卜长\Person{塔革斯}，或凭猜测或其自己的观察，说：\Quote{这里有亵圣的人，阻碍了仪式。}于是\Person{戴克里先}勃然大怒，下令不仅所有参与圣礼的人，而且所有住在宫内的人，都必须献祭，若拒绝则受鞭笞。并且，他通过给指挥官的诏书，命令所有士兵都须被迫行此亵渎之事，违者开除军籍。他的愤怒至此为止；但在那个季节，他没有对天主的法律和宗教再做更多的事。过了一些时候，他前往\Place{比提尼亚}过冬；不久\Person{伽列里乌斯·凯撒}也到了那里，满怀狂怒的怨恨，企图煽动那位轻率的老人进行他已开始的对基督徒的迫害。我得知他愤怒的原因如下。\par

% structure: p0044 source=h2 target=section reason=relative-heading-rank; base=h2

\section{第十一章}

\Person{伽勒里乌斯}的母亲，一位极其迷信的妇人，是山神的崇拜者。她性情如此，几乎每日献祭，并用祭过偶像的肉宴请她的仆人；但她家中的基督徒不肯参加这些宴席；当她与外邦人同宴时，他们却持守斋戒和祈祷。因此，她对基督徒心怀恶意，并以妇人的怨言煽动她那与她同样迷信的儿子去消灭他们。于是，整个冬天，\Person{戴克里先}与\Person{伽勒里乌斯}单独举行会议，无他人参与；众人皆认为他们所商讨的是帝国最重大的事务。那位老人长期抵制\Person{伽勒里乌斯}的狂怒，指出在全世界引发骚乱、流那么多血是何等有害；基督徒惯于欣然赴死；他只需将此种宗教的人逐出宫廷和军队便已足够。然而，他无法抑制那固执之人的疯狂。于是，他决定征求朋友们的意见。\Person{戴克里先}的性格中有这样一个特点：当他决定行善时，他不经商议便行，以便所有赞美归于自己；但当他决定作恶——他知道这会被指责——他便召来许多顾问，把自己的过错归咎于他人。因此，几位民政长官和几位军事将领被准许提供意见；问题按官阶高低依次提出。有些人对基督徒怀有私怨，认为应将其铲除，视其为众神的敌人和既定宗教仪式的反对者。另一些人持有不同看法，但了解了\Person{伽勒里乌斯}的意愿后，或出于惧怕得罪，或出于取悦之心，都附合了反对基督徒的意见。然而，即便如此，皇帝仍未被说服。他决意首先咨询他的众神；为此，他派遣一名占卜者前往\Place{米利都}的\Person{阿波罗}神谕所询问，得到的回答正如神圣宗教的敌人所料。于是\Person{戴克里先}被说服了。虽然他不能再违抗他的朋友、\Person{凯撒}和\Person{阿波罗}，但他仍试图保持节制，命令此事在不流血的情况下进行；而\Person{伽勒里乌斯}则希望将所有拒绝献祭的人活活烧死。\par

% structure: p0046 source=h2 target=section reason=relative-heading-rank; base=h2

\section{第十二章}

为完成此事，选了一个合适而吉利的日子；\OriginalTerm{堤尔米努斯}{Terminus}神的节日（在三月初七）被优先选出，仿佛是要借此终止基督徒的宗教。\par

\Quote{那日，死亡的前兆升起，\newline{}万恶之始，长久苦难的根源；}\par

这苦难不仅降临在基督徒身上，也降临在整个大地。那天破晓时分，正值\Person{戴克里先}第八次执政、\Person{马克西米安}第七次执政之际，天刚蒙蒙亮，长官、将领、护民官和国库官员突然来到\Place{尼科米底亚}的教堂，强行破门后，到处搜寻神像。找到了圣经圣书，便投入火中；教堂的器具和家具被洗劫一空：到处是掠夺、混乱、喧嚣。那座教堂建在高地上，从王宫可以望见；\Person{戴克里先}和\Person{伽勒里乌斯}站在上面如同瞭望塔，长久争论是否该放火。\Person{戴克里先}的意见占了上风，他担心一旦大火燃起，城市的一部分可能被烧毁，因为教堂周围有许多高大的建筑。于是，近卫军全副武装，带着斧头和其他铁器，四处冲杀，几个小时之内便将那座极高的建筑夷为平地。\par

% structure: p0050 source=h2 target=section reason=relative-heading-rank; base=h2

\section{第十三章}

次日，一道敕令颁布，剥夺基督徒的一切荣誉和官职；并规定，无论其地位高低，都应受酷刑，任何诉讼都可针对他们提起；而另一方面，他们不得在涉及伤害、通奸或偷窃的案件中作为原告；最终，他们既无获得自由的能力，也无选举权。有一个人撕下这道敕令，并将其撕碎，虽然行为不当，但精神高尚，并轻蔑地说：\Quote{这些是哥特人和萨尔马提亚人的胜利。}他立刻被抓住并受审，不但被施以酷刑，还依法被活活烧死；他表现出令人钦佩的忍耐，最终化为灰烬。\par

% structure: p0052 source=h2 target=section reason=relative-heading-rank; base=h2

\section{第十四章}

但\Person{加勒里乌斯}不满足于敕令的措辞，另寻方法以左右皇帝。为促使他在迫害中施以过度的残忍，他雇佣私人密使纵火烧宫；宫殿一部被焚，罪责便落在基督徒身上，视其为公敌；因这场火灾，基督徒的名号变得可憎。据说基督徒与宦官合谋，意图毁灭君王；两位君王几乎在自己宫中活活烧死。\Person{戴克里先}，总是自诩精明睿智，却对这场阴谋毫无察觉，反被怒气点燃，立即下令严刑拷打他所有家奴，逼其招供阴谋。他坐于审判席，眼见无辜之人在火刑中受折磨以求真相。所有官员及宫廷总管皆获特别授权施以酷刑；众人争相竞先，看谁先揭露这桩阴谋。然而，事实真相始终未在任何地方被发现；因为无人对\Person{加勒里乌斯}的家奴施用酷刑。\Person{加勒里乌斯}本人始终伴随\Person{戴克里先}，不断催促，不让这位轻率老人的怒火冷却。随后，相隔十五天后，他又试图第二次纵火；但火势迅速被察觉并扑灭。然而，纵火者仍未查明。就在那一天，正值仲冬准备离开的\Person{加勒里乌斯}突然匆忙出城，宣称他是为了逃命免被烧死。\par

% structure: p0054 source=h2 target=section reason=relative-heading-rank; base=h2

\section{第十五章}

于是\Person{戴克里先}大怒，不仅对自家家奴，而且无差别地向所有人施暴；他首先强迫女儿\Person{瓦莱里亚}和妻子\Person{普里斯卡}被祭献玷污。曾一度权势熏天、在宫廷和皇帝面前拥有最高权威的宦官被处死。司铎和教会其他官员被捕，没有证人证据或口供，便被定罪，连同家人一同被处决。活活烧死时，不分性别年龄；因人数众多，他们不被逐一烧死，而是一群人同受一圈火焰；仆人脖子上绑着磨石，被投入海中。对天主其余子民的迫害也同样严酷；因为法官分散到各神庙，试图强迫每个人祭献。监狱爆满；前所未闻的酷刑被发明；为防止基督徒无意中得到公正审判，在法院中、审判席旁放置祭台，每个诉讼人在案件审理前必须先焚香。因此，人们向法官求告就如同向神明求告一般。命令也下达给\Person{马克西米安·赫库利乌斯}和\Person{君士坦提乌斯}，要求他们同意执行敕令；因为在如此重大的事务上，从未征求过他们的意见。\Person{赫库利乌斯}性情不仁，欣然服从，并在他的\Place{意大利}领土上强制执行敕令。\Person{君士坦提乌斯}则不然，为避免显得违抗上级命令，他允许拆毁教堂——不过是些墙壁，可以重建——但他完好保全了天主真正的圣殿，即人的身体。\par

% structure: p0056 source=h2 target=section reason=relative-heading-rank; base=h2

\section{第十六章}

于是全地受苦；从东到西，除\Place{高卢}地区外，三只贪婪的野兽继续肆虐。\par

\Quote{就算我有一百张嘴、一百条舌头，\newline{}铜嗓子、铁肺，\newline{}也无法描述那可怕场面的一半。}\par

也无法详述各地统治者加诸虔诚无辜之人的惩罚。\par

但有何必要逐一细述这些事，尤其对你，我最亲爱的\Person{多纳图斯}，你比众人更暴露在那场猛烈迫害的暴风雨中？因为你曾落入总督\Person{弗拉基尼安}——一个不小的凶手——之手，随后落入\Person{希罗克勒斯}之手，他由副职升任\Place{比提尼亚}总督，是迫害的发起者和顾问，最后落入其继任者\Person{普里斯基利安}之手，你向人类展示了不可战胜的宽宏气度。你九次受刑架和各种折磨，九次以荣耀的信德宣认挫败仇敌；在九次战斗中你制胜魔鬼及其精兵；以九次胜利战胜这世界及其恐怖。天主看见你作为胜利者，驾着战车，拖着的不是白马，也不是巨象，而是那些曾俘虏万民的人，这景象何其悦乐！如此凌驾于世上的主宰，才是真正的凯旋！如今，因你的英勇，他们被征服了，当你藐视他们邪恶的敕令，并藉着坚定的信德和灵魂的刚毅，溃败了暴政权威一切虚妄的恐怖。鞭子、铁爪、火、剑、各种酷刑都对你无效；任何暴力都无法夺去你的忠诚和持守的决心。这就是作为天主的门徒，作为基督的士兵；这士兵，没有敌人能将他从天上的营寨中赶出，没有豺狼能将他夺去；没有诡计能陷害他，没有身体的痛苦能制伏他，没有折磨能推翻他。最后，经过那九场光荣的战斗，魔鬼被你击败，他不敢再与你争战，因为他屡次试探后发现你不可战胜；他不再挑战你，恐怕你已经抓住了那伸向你的胜利花冠；一个永不凋谢的花冠，虽然你尚未领取，但已因你的德行和功劳安放在主的国度里。但现在让我们回到叙述的进程。\par

% structure: p0061 source=h2 target=section reason=relative-heading-rank; base=h2

\section{第十七章}

邪恶的计划既已付诸实施，戴克里先如今已被幸运所弃，立即动身前往罗马，在那里庆祝他统治第二十年的开端。那典礼于十二月朔日前第十二日举行；忽然间，皇帝无法忍受罗马人的言论自由，愠怒而急躁地离城而去。一月朔日临近，那天将第九次授予他执政官之职；然而，他不愿在罗马多待十三天，宁可选择在拉韦纳首次以执政官身份出现。然而，他在严冬中启程，正值酷寒与连绵阴雨，染上了一种轻微却缠绵的疾病；这病不断折磨他，以至于他大部分时间不得不由人抬在轿中。随后，在夏末之际，他沿着多瑙河岸绕行，如此来到了尼科米底亚。他的病情此时已变得更加严重和压抑；但他仍然命人将自己抬出，以便奉献那座他在统治第二十年结束时建造的竞技场。他立刻变得如此虚弱无力，以致向所有神明祈求他的生命。然后，在十二月望日，宫中突然传来悲伤、哭泣和哀悼之声，朝臣们四处奔忙；全城一片寂静，传言戴克里先已死，甚至已下葬；但次日清晨，又突然谣传他还活着。此时，他的仆从和朝臣的脸色由忧郁转为欢快。然而，仍有人怀疑他的死被隐瞒，直到加莱里乌斯·凯撒到来，唯恐在此期间军队企图改变政体；这一猜疑变得如此普遍，以至于没有人相信皇帝还活着，直到三月朔日，他公开露面，但面容如此苍白（他的病已持续近一年），几乎认不出来。那类似死亡的昏厥发生在十二月望日；尽管他有所恢复，却再未达到完全健康，因为他精神变得错乱，有时疯癫，有时清醒。\par

% structure: p0063 source=h2 target=section reason=relative-heading-rank; base=h2

\section{第十八章。}

几天之内，加莱里乌斯·凯撒抵达，并非祝贺其岳父恢复健康，而是迫使他逊位。他早已催促马克西米安·赫尔克里乌斯同样行事，并以内战的恐惧使这位老人屈服；如今他攻击戴克里先。起初，他以温和友善的语气，说年事已高且日益衰弱的身体使戴克里先无法承担治国之责，他需要在辛劳之后稍事休息。加莱里乌斯为证实其论点，举出内尔瓦的例子，他将帝国的重担交给了图拉真。\par

但戴克里先回答说，一个曾居于超乎众人、显赫地位的人，降入卑微的境地是不合适的；而且也不安全，因为在如此漫长的统治中，他必然树敌众多。内尔瓦的情况则大为不同：他仅在位一年，便感到自己或因年迈或因不谙政务，无法胜任如此重大的事务，因而抛下政府的舵柄，回到他早已在其中度过晚年的私人生活。但戴克里先补充道，如果加莱里乌斯想要皇帝的称号，没有什么能阻止将之授予他和君士坦提乌斯，以及马克西米安·赫尔克里乌斯。\par

加莱里乌斯的想象早已攫取了整个帝国，看出这一提议除了一个空洞的名称之外，几乎不会给他带来什么，于是回答说，戴克里先本人所作的规定应不可侵犯；该规定确立了应有两位较高级别者拥有最高权力，另有两位较低级别者予以辅佐。两个平等者之间容易保持和谐，但四个则永不可能；他说，如果戴克里先不逊位，他就必须为自己的利益着想，以免继续处于较低级别，且是该级别的末位；过去十五年来，他如同流放一般被局限在伊利里亚和多瑙河沿岸，不断与野蛮民族搏斗，而其他人则安逸地统治着比他更广阔、更文明的领土。\par

戴克里先早已通过马克西米安·赫尔克里乌斯的来信，得知加莱里乌斯在他们会谈时所说的一切，也得知他正在增兵；此时听到他的话语，这位毫无气概的老人流下眼泪，说道：\Quote{就照你的意思吧。}\par

剩下的是要共同推选凯撒。\Quote{可是，}加莱里乌斯说，\Quote{何必征求马克西米安和君士坦提乌斯的意见？因为他们对我们所做的一切必然顺从。}——\Quote{他们当然会，}戴克里先答道，\Quote{因为我们得选举他们的儿子。}\par

马克西米安·赫尔克里乌斯有一个儿子，名叫马克森提乌斯，娶了加莱里乌斯的女儿为妻；此人品性邪恶而有害，且极其傲慢固执，从不向父亲或岳父行应有的敬礼，因此为他们两人所憎恶。君士坦提乌斯也有一个儿子，名叫君士坦丁，是一位非常卓越的青年，完全配得上凯撒的高位。他那出众的英俊仪表、对一切军务的严格专注、端正的品行和非凡的和蔼可亲，使他深受军队爱戴，成为每个人的选择。他当时在宫中，早已被戴克里先任命为第一等护民官。\par

\Quote{该当如何？} \Person{加勒里乌斯}说，\Quote{马克森提乌斯不配此职。他尚为平民时即辱我，一旦得权，又将如何？}—— \Quote{但君士坦丁和蔼可亲，若由他治理，今后在世人眼中，必超越其父之温和美德。}—— \Quote{若我的意愿与判断不被重视，那就如此吧。应任命那些听命于我、畏惧我、非我命令不敢行事之人。}—— \Quote{那么该任命谁？}—— \Quote{塞维鲁。}—— \Quote{怎可！那个舞者，那个惯常醉汉，黑夜当作白昼，白昼当作黑夜？}—— \Quote{他当此职，因他证明自己是忠实的军需官与军粮供给者；并且，我已派他往\Person{马克西米安}处领取紫袍。}—— \Quote{好，我同意；但你还建议谁？}—— \Quote{他，}\Person{加勒里乌斯}指着\Person{达亚}——一个半蛮族的青年——说道。原来\Person{加勒里乌斯}最近将自己名字的一部分赐予这青年，称他为\Person{马克西米努斯}，如同昔日\Person{戴克里先}将\Person{马克西米安}之名赐予\Person{加勒里乌斯}，为取兆头，因\Person{马克西米安·赫库利乌斯}曾以坚定忠信侍奉他。—— \Quote{你举荐的是谁？}—— \Quote{我的亲戚。}—— \Quote{唉！}\Person{戴克里先}深叹一口气，\Quote{你所推荐之人，并非适合管理公务之材！}—— \Quote{我已考验过他们。}—— \Quote{那么，既然你即将执掌帝国统治，就由你来负责吧；至于我，当我继续做皇帝时，我辛勤劳苦，长久谨慎，为保障国家安康；如今若有任何灾祸发生，责不在我。}\par

% structure: p0071 source=h2 target=section reason=relative-heading-rank; base=h2

\section{第十九章}

事情如此商定后，\Person{戴克里先}和\Person{加勒里乌斯}列队出发，公布凯撒的任命。所有人都看着\Person{君士坦丁}；因为毫无疑问，人选会落在他身上。在场的军队，以及其他被召集参加典礼的军团的主要士兵，都将目光投向\Person{君士坦丁}，为他即将当选而欢欣鼓舞，并专心祈祷他兴旺发达。离\Place{尼科米底亚}约三里处有一高地，\Person{加勒里乌斯}曾在那里接受紫袍；那里立有一根柱子，上有\Person{朱庇特}雕像。队伍前往那里。召集士兵集会。\Person{戴克里先}含泪向他们演说，说他身体衰弱，疲惫后需要休息，他将把帝国交给更有活力与能力的人，同时任命新的凯撒。观众极其热切地等待任命。突然他宣布凯撒是\Person{塞维鲁}和\Person{马克西米努斯}。众人无不惊愕。\Person{君士坦丁}站在公众可见之处，人们开始彼此询问他的名字是否也被改为\Person{马克西米努斯}；此时，在众人注视下，\Person{加勒里乌斯}向后伸手，将\Person{君士坦丁}推开，把\Person{达亚}拉到前面，脱去他的平民服装，把他放在最显眼的位置。所有人都诧异他是谁，从何而来；但无人敢于干涉或提出异议，他们的心思被这奇怪而意外的事件搞糊涂了。\Person{戴克里先}脱下紫袍，披在\Person{达亚}身上，并恢复自己原来的名字\Person{狄奥克勒斯}。他走下讲台，乘车穿过\Place{尼科米底亚}；然后这位老皇帝，如同一位退伍的士兵从军役中解脱，被送归故里；而\Person{达亚}，不久前才从森林中放牧牛羊被征召为普通士兵，立刻成为近卫军，旋即升为护民官，次日成为凯撒，获得了践踏和压迫东方帝国的权力；一个既不懂战争也不懂民事的人，从牧人变成了军队统帅。\par

% structure: p0073 source=h2 target=section reason=relative-heading-rank; base=h2

\section{第二十章}

\Person{加勒里乌斯}实现了驱逐两位老人后，开始自视为罗马帝国的唯一君主。任命\Person{君士坦提乌斯}为首位是出于必要；但\Person{加勒里乌斯}并不看重这个性情温和、健康衰退且不稳定的人。他预料\Person{君士坦提乌斯}会很快死去。即使这位君主康复，似乎也不难迫使他脱下皇袍；否则，如果三位同僚联合逼他退位，他还能做什么？\Person{加勒里乌斯}始终将\Person{李锡尼}带在身边，\Person{李锡尼}是他的老朋友，也是最早的战友，他咨询\Person{李锡尼}处理所有事务；但他不愿以儿子的称号任命\Person{李锡尼}为凯撒，因为他打算在\Person{李锡尼}取代\Person{君士坦提乌斯}后，以兄弟的称号任命他为皇帝，而他自己则掌握主权，以无限自由统治全球。此后，他打算举行二十周年庆典；然后将凯撒之职授予他当时九岁的儿子\Person{坎迪狄安}；最后，像\Person{戴克里先}那样退位。这样，\Person{李锡尼}和\Person{塞维鲁}为皇帝，\Person{马克西米努斯}和\Person{坎迪狄安}为下一等凯撒；他幻想自己仿佛被一道无法攻破的城墙环绕，将在安全与和平中度过晚年。这就是他的计划；但被他视为对手的天主挫败了所有这些空想。\par

% structure: p0075 source=h2 target=section reason=relative-heading-rank; base=h2

\section{第二十一章}

他既如此获得了最高权力，便一心要折磨那个他为自己开辟道路的帝国。波斯人的习俗和惯例是民众自身成为国王的奴隶，而国王则视民众为奴隶。这个邪恶的人，自从战胜波斯人以后，便不知羞耻地不断赞扬这种制度，并决心在罗马领地内加以确立；由于无法通过明文法律实现，他便仿效波斯国王的行径，剥夺人们的自由。他首先贬低那些他打算惩罚的人；然后不仅下级官员被他施以酷刑，就连城市中的首领和地位最显赫的人也不例外，而且这些惩罚往往涉及细微之事和民事问题。钉十字架是死刑的即时惩罚；较轻的罪行则用枷锁。尊贵的贵妇被拖入劳役所；当有人要受鞭打时，便在地上固定四根木桩，将他绑在上面，这是奴隶刑罚中从未见过的方式。我该说他那娱乐的厅堂和他喜爱的消遣吗？他养了许多熊，这些熊在凶猛和体魄上最像他自己，是他统治期间收集来的。每当他想要放纵性情时，便下令带进某只熊，把人扔给那野兽，与其说是被吞噬，不如说是被吞没；当他们的肢体被撕裂时，他极其满意地大笑；他从未在不观看流血的情况下用晚餐。平民被判活活烧死；他通过针对基督徒的法令开始了这种处决方式，命令在酷刑和定罪之后，用慢火焚烧他们。他们被绑在柱子上，先用适度的火焰灼烧脚底，直到肌肉因燃烧而收缩，从骨头上撕裂；然后点燃又熄灭的火把指向他们身体的各个肢体，使任何部分都不得幸免。同时，不断往他们脸上泼冷水，润湿他们的嘴巴，以免因下颌干燥而死。最后，当数小时后，炽热已烧毁皮肤并渗入内脏时，他们才死去。死尸被放在火葬堆上完全烧毁；他们的骨头被收集起来，磨成粉末，扔进河或海里。\par

% structure: p0077 source=h2 target=section reason=relative-heading-rank; base=h2

\section{第二十二章}

如今，他在折磨基督徒时学会的那种残忍已成为习惯，并毫无区别地对所有人施行。他不习惯施加较轻的惩罚，如放逐、监禁或送罪犯去矿山劳作；而是焚烧、钉十字架、投喂野兽，这些事每天发生，毫不犹豫。对于较小的过错，他自家的仆人和管家被用长矛而不是棍棒责打；对于大罪，斩首是极少数人能得到的优待；而且那似乎是一种恩惠，因旧日的效力，才允许人以最轻松的方式死去。但这些与接下来相比，只是\Person{加勒里乌斯}统治中的小恶。因为雄辩被消灭，辩护人被除掉，精通法律的人或被流放或被杀害。有用的书籍被视同巫术和禁术；所有拥有它们的人都遭到践踏和咒骂，仿佛他们是政府的敌人和公敌。法律解体了，法官被赋予无限制的自由——这些法官从士兵中选出，粗鲁无文，被放任到各省，没有辅佐官来指导或控制他们。\par

% structure: p0079 source=h2 target=section reason=relative-heading-rank; base=h2

\section{第二十三章}

然而，引起公共和普遍灾难的，是同时向各省和各城市征收的税。勘测员被派往各处，进行普遍而严厉的审查，可怕的景象呈现出来，如同胜利的敌人的暴行和俘虏的悲惨状况。每一块土地都被测量，葡萄树和果树被清点，各种动物被登记，并编造了人头税册。城市中，平民无论居住在城墙内还是城墙外，都被集合起来，市场挤满了成群的家族，都带着他们的孩子和奴隶，拷打和鞭笞的声音回荡，儿子被挂在刑架上以逼供父亲的财产，最忠实的奴隶被迫在痛苦中作证反对他们的主人，妻子作证反对她们的丈夫。在缺乏其他证据时，人们被拷打以让他们自己指控自己；一旦痛苦迫使他们承认子虚乌有之事，那些想象中的财产就被记在册上。无论是青年、老年还是疾病，都得不到任何豁免。病人和体弱者被抬来；每个人的年龄被估算；为了增加人头税，给年轻人加岁数，从老年人那里减岁数。普遍的哀恸和悲伤弥漫。按照战争法，征服者对被征服者所做的一切，此人竟胆敢对罗马人和罗马的臣民施做，因为他的祖先曾因胜利的\Person{图拉真}所征收的类似税款而负有责任，这是对达契亚人频繁叛乱的惩罚。此后，按人头征收钱款，仿佛是为生存自由支付了代价；然而，他们并不完全信任同一批勘测员，而是一批又一批地被派去进行进一步查探；于是税款加倍，不是因为新勘测员有了新的发现，而是因为他们随心所欲地增加了前一次的税额，以免显得他们被雇来一无所用。与此同时，动物数量减少，人死去；然而，甚至死者也要纳税，以至于没有人能活着或死去而不受征税。只剩下乞丐，从他们那里什么都勒索不到，他们的贫穷和悲惨反倒保护了他们免受虐待。但是这位虔敬的人对他们动了怜悯之心，决定不再让他们生活在贫困中，于是他命人将他们全部集合起来，装上船，然后沉入海中。他是如此仁慈，安排在他的治理下无人会匮乏！就这样，当他采取有效措施，使任何人都不能以假装的贫穷为借口逃税时，他杀害了许多真正的可怜人，违背了所有人道法则。\par

% structure: p0081 source=h2 target=section reason=relative-heading-rank; base=h2

\section{\Emph{第二十四章。}}

天主的审判已经临近他，随之而来的是他的命运开始衰败和凋零的季节。当他正以我上述的方式忙碌时，他并没有着手推翻或驱逐\Person{君士坦提乌斯}，而是等待他的死亡，但并没有想到死亡会来得如此之快。\Person{君士坦提乌斯}病得很重，写信给\Person{伽列里乌斯}，请求允许他的儿子\Person{君士坦丁}来见他。他早已提出过同样的请求，但徒劳无功；因为\Person{伽列里乌斯}根本不打算答应。相反，他对那年轻人屡次设下圈套，因为他不敢公开使用暴力，以免激起内战，招致他最惧怕的——军队的仇恨和愤怒。以男子汉的训练和娱乐为借口，他让\Person{君士坦丁}与野兽搏斗；但这一计谋失败了，因为天主的全能保护了\Person{君士坦丁}，并在危难时刻将他从\Person{伽列里乌斯}手中解救出来。最后，\Person{伽列里乌斯}再也无法回避\Person{君士坦提乌斯}的请求，于是在一天晚上给了\Person{君士坦丁}一份出发的许可，命令他次日早晨带着皇帝的信件出发。\Person{伽列里乌斯}打算要么找借口扣留\Person{君士坦丁}，要么命令\Person{塞维鲁}在路上逮捕他。\Person{君士坦丁}看穿了他的意图，因此，晚餐后，当皇帝就寝时，他匆忙离开，从主要驿站带走了所有公费供养的马匹，然后逃脱了。次日，皇帝故意在卧室里待到中午，然后命令召见\Person{君士坦丁}；但他得知\Person{君士坦丁}在晚餐后就出发了。他怒火中烧，命令备好马匹去追捕并拖回\Person{君士坦丁}；当听说大路上所有的马都被带走了，他几乎忍不住落泪。与此同时，\Person{君士坦丁}以不可思议的速度行进，到达了他父亲那里，父亲已奄奄一息。\Person{君士坦提乌斯}将儿子托付给士兵，将最高权力交到他手中，然后如他所愿，安详平静地去世。\par

\Person{君士坦丁皇帝}即位后，将恢复基督徒敬拜的活动和对他们天主的敬拜作为他的首要之事；于是他以恢复圣教开始了他的治理。\par

% structure: p0084 source=h2 target=section reason=relative-heading-rank; base=h2

\section{\Emph{第二十五章。}}

几天之后，\Person{君士坦丁}的肖像饰以月桂，被送到那有害的野兽面前，好让他通过接受这标记，承认\Person{君士坦丁}为皇帝。他犹豫良久，不知是否该接受，几乎要将肖像和送肖像者一并投入火中，但他的亲信劝阻了这疯狂的决议。他们告诫他危险所在，并说明：如果\Person{君士坦丁}率军前来，所有士兵——那些违背其意愿而被立为默默无闻或不知名的凯撒的人——都会承认他，并争先恐后地投奔他的旗帜。于是\Person{加莱里乌斯}，尽管极不情愿，接受了肖像，并将皇帝紫袍送给\Person{君士坦丁}，好让他显得是自愿接纳这位王子与他共享权力。如今他的计划被打乱了，他不能再像先前打算的那样，在不超出皇帝限定数目的情况下接纳\Person{李锡尼}。但他想出了这个办法：让较为年长的\Person{塞维鲁}被立为皇帝，而\Person{君士坦丁}，不授予他已被任命的皇帝称号，而是与\Person{马克西米努斯·达亚}共同接受凯撒称号，从而从第二位降为第四位。\par

% structure: p0086 source=h2 target=section reason=relative-heading-rank; base=h2

\section{第二十六章}

事情在某种程度上似乎已按\Person{加莱里乌斯}的意愿安排妥当，这时又传来一个警报：他的女婿\Person{马克森提乌斯}在\Place{罗马}被立为皇帝。原因如下：\Person{加莱里乌斯}决意通过永久性的税收来吞噬帝国，竟狂妄到连\Place{罗马}人民也不免除这种奴役。于是，收税官被派往\Place{罗马}，开列公民名单。大约同时，\Person{加莱里乌斯}裁减了禁卫军。\Place{罗马}还剩有一些该军团士兵，他们趁机处死了一些地方官员，并在喧闹的民众默认下，给\Person{马克森提乌斯}穿上了皇帝紫袍。\Person{加莱里乌斯}听到这个消息，虽为这突然之事感到不安，但并未十分惊慌。他憎恨\Person{马克森提乌斯}，而且他不能再将凯撒的尊号赐予他，因为已有两人（\Person{达亚}和\Person{君士坦丁}）享有此尊号；此外，他认为自己曾不情愿地赐予那尊号一次已足够了。于是他召来\Person{塞维鲁}，劝勉他重新夺取自己的统治权和主权，并命他统率\Person{马克西米安·赫库利乌斯}原先指挥的那支军队，去进攻\Place{罗马}的\Person{马克森提乌斯}。在那里，\Person{马克西米安}的士兵曾多次受到各种奢侈款待，因此他们不仅有兴趣保护这座城市，而且渴望定居在那里。\par

\Person{马克森提乌斯}深知自己罪孽深重；尽管他对父亲军队的效劳有着世袭的权利，并可能希望将其争取过来，但他想到这也会引起\Person{加莱里乌斯}的同样考虑，可能导致\Person{加莱里乌斯}将\Person{塞维鲁}留在\Place{伊利里库姆}，亲自率军进攻\Place{罗马}。在这种忧虑下，\Person{马克森提乌斯}寻求保护自己免受悬在头上的危险。他派人将紫袍送给退位后居住在\Place{坎帕尼亚}的父亲，并再次尊他为奥古斯都。善变的\Person{马克西米安}急切地重新穿上他不情愿脱下的紫袍。同时，\Person{塞维鲁}进军，率军逼近\Place{罗马}城墙。很快，士兵们举起军旗，抛弃\Person{塞维鲁}，归顺了他们前来讨伐的\Person{马克森提乌斯}。被抛弃的\Person{塞维鲁}除了逃跑还有什么办法？他遇上了已经恢复帝位的\Person{马克西米安}。于是他逃到\Place{拉文纳}，带着少数士兵闭城固守。但他意识到自己将被交出去，便自愿投降，将紫袍归还给他曾接受紫袍的人；此后他得到的唯一恩惠是轻松的死亡：他被迫割开血管，就这样温和地死去了。\par

% structure: p0089 source=h2 target=section reason=relative-heading-rank; base=h2

\section{第二十七章}

但是\Person{马克西米安}深知\Person{加莱里乌斯}的暴烈性情，开始思虑，\Person{加莱里乌斯}闻悉\Person{塞维鲁}的死讯后必会怒火中烧，进军\Place{意大利}，并且可能得到\Person{达亚}的协助，从而调集过于强大的兵力而无法抵御。因此，在加固了\Place{罗马}城，并为防御战作了周密准备之后，\Person{马克西米安}前往\Place{高卢}，以便将他年幼的女儿\Person{福斯塔}嫁给\Person{君士坦丁}，从而赢得那位王子的支持。与此同时，\Person{加莱里乌斯}集结军队，侵入\Place{意大利}，并向\Place{罗马}推进，决心消灭元老院，将全体人民斩于剑下。但他发现一切都被封锁并加固了，无法以强攻夺取，而围城又是一件艰巨的任务；因为\Person{加莱里乌斯}没有带来足以包围城墙的军队。也许由于从未见过\Place{罗马}，他想象它只比他熟悉的那些城市大不了多少。但他的几个军团厌弃了一个父亲对他女婿的邪恶企图，以及罗马人对罗马的邪恶行动，遂放弃了他的权威，将他们的军旗投向了敌人。他剩余士兵已经开始动摇，这时\Person{加莱里乌斯}害怕遭遇与\Person{塞维鲁}相似的命运，他傲慢的精神被击碎并屈辱，跪倒在他的士兵脚下，不断恳求他们不要把他交给敌人，直到他许下重赏才说服了他们。于是他撤离\Place{罗马}，溃不成军地逃窜。若有人哪怕率领一小队人马追击，他很容易在逃亡中被截断。他意识到危险，便允许士兵四散，并广泛地劫掠和破坏，以便若有追兵，就会在一片毁坏之地缺乏一切生存资料。于是，这支瘟疫般的队伍所经之处，\Place{意大利}各地均被蹂躏，一切被洗劫，主妇被强暴，处女被侵犯，父母和丈夫在酷刑下被迫交代他们藏匿财物、妻子和女儿的地方；成群的牲畜像从蛮族那里掠得的战利品一样被赶走。就这样，他曾经是罗马皇帝，如今却成了\Place{意大利}的蹂躏者，在使所有人无差别地遭受战争灾难之后，退入自己的领地。确实，早在他获得最高权力的那个时候，他就已自称为罗马之名的敌人；他曾提议帝国应称为达西亚帝国，而非罗马帝国。\par

% structure: p0091 source=h2 target=section reason=relative-heading-rank; base=h2

\section{第二十八章。}

\Person{加莱里乌斯}逃走后，\Person{马克西米安}从\Place{高卢}返回，与他的儿子共同执掌大权；但人们对年轻人的服从多于对年长者：因为\Person{马克森提乌斯}权力最大，在位时间最长；而且在此次事件中，\Person{马克西米安}正是仰赖他才得以恢复帝位。老人不甘被剥夺不受约束的统治权，并以幼稚的竞争心嫉妒自己的儿子；因此他开始谋划如何驱逐\Person{马克森提乌斯}，恢复他昔日的统治。这看似容易，因为那些背叛\Person{塞维鲁}的士兵原本是在他自己的军队中服役的。他召集了\Place{罗马}民众和士兵大会，仿佛要就公共事务的悲惨状况发表演说。在就此主题说了许多之后，他伸出手指向他的儿子，指控他是一切祸患的制造者和国家灾难的首要原因，然后撕下了他肩上的紫袍。被如此剥去衣袍的\Person{马克森提乌斯}从讲坛上一跃而下，被士兵们接在臂膀中。他们的愤怒和喧嚣使那个不自然的老人陷于混乱，于是他像另一个高傲者\Person{塔克文}一样，被赶出了\Place{罗马}。\par

% structure: p0093 source=h2 target=section reason=relative-heading-rank; base=h2

\section{第二十九章。}

随后\Person{马克西米安}回到\Place{高卢}；在那里停留一段时间后，他去了他儿子的仇敌\Person{加莱里乌斯}那里，假装就公共秩序的安定进行商议；但他的真正目的乃是假借和解之名，寻找谋杀\Person{加莱里乌斯}的机会，夺取他的那部分帝国，以替代他自己已被各地排斥的份额。\par

\Person{戴克里先}正在\Person{加莱里乌斯}的宫廷，当时\Person{马克西米安}到达；因为\Person{加莱里乌斯}打算现在授予\Person{李锡尼}以代替\Person{塞维鲁}的最高权力标记，最近已将\Person{戴克里先}召来参加盛典。于是在他和\Person{马克西米安}面前举行了仪式；这样，同一时间有六人统治着帝国。\par

\Person{马克西米安}的计谋受挫后，他像此前两次一样逃回\Place{高卢}，心怀恶念，企图用诡计欺骗\Person{君士坦丁}——此人不仅是他的女婿，更是他儿子的孩子。为更易得手，他卸下了皇室的紫袍。法兰克人已举兵。\Person{马克西米安}建议毫无戒心的\Person{君士坦丁}不要率全军前往，并说只需少数士兵便足以制服那些蛮族。他提出此建议，为的是留下军队供自己收买，并使\Person{君士坦丁}因兵力薄弱而被击败。少年君主因这建议出自一位年迈而有经验的统帅，便信以为明智；又因出自岳父之口，便听从了。他出发了，将大部分军队留在后方。\Person{马克西米安}等待数日，一旦估算\Person{君士坦丁}已进入蛮族境内，便突然重披紫袍，夺取国库，照例向士兵大发赏赐，并谎称\Person{君士坦丁}遭遇了不久后降临在他自己身上的灾难。\Person{君士坦丁}立即得知此事，以惊人的速度率军赶回。\Person{马克西米安}尚未做好迎战准备，被出其不意地制服，士兵们回归本职。\Person{马克西米安}占据了\Place{马赛}（他逃往那里），关闭城门。\Person{君士坦丁}走近，见\Person{马克西米安}在城墙上，便以并无严厉或敌意的言辞向他说话，询问他的意图、他想做什么、为何做出如此特别不体面的行为。但\Person{马克西米安}从城墙上不停地辱骂和诅咒\Person{君士坦丁}。忽然，对面的城门被打开，围城者进入城中。那位叛乱的皇帝、不慈的父亲、背信的岳父被拖到\Person{君士坦丁}面前，听取了对他的罪行的陈述，被剥去皇袍，受斥责后保全了性命。\par

% structure: p0097 source=h2 target=section reason=relative-heading-rank; base=h2

\section{第三十章。}

\Person{马克西米安}既如此丧失了作为皇帝和岳父应有的尊敬，便对自己的卑下处境日益不耐，且因未受惩罚而胆大妄为，又对\Person{君士坦丁}设下新的阴谋。他去找女儿\Person{福斯塔}，既用恳求，又用谄媚的哄骗，求她背叛丈夫。他许诺为她寻得比与\Person{君士坦丁}的婚姻更荣耀的联姻，并让她打开皇帝的卧室，不予设防。\Person{福斯塔}答应做他所求的一切，并立刻将全部事情告知了丈夫。他们设下计谋，要在\Person{马克西米安}正行凶时当场拿获。他们安排一个卑贱的宦官代替皇帝被杀害。夜深时分，\Person{马克西米安}起身，见诸事利于他的奸计。守卫的士兵很少，且离卧室颇远。然而，为避免猜疑，他上前对他们说，他做了一个梦，想告诉他的女婿。他持械进入，杀死宦官，得意地跃出，承认了谋杀。正在此时，\Person{君士坦丁}从对面带着一队士兵出现；尸体被从卧室抬出；凶手持凶当场被拿，惊愕万分，\par

\Quote{如石站立，默然不动；}\par

而\Person{君士坦丁}斥责他的不敬和滔天罪行。最后，\Person{马克西米安}获准自择死法，便自缢而亡。\par

那位最强大的\Place{罗马}君主——曾如此长久地以极其荣耀统治，并庆祝了他的二十周年纪念——那位最傲慢的人，就这样断了颈骨，以可耻卑贱的死亡结束了他可憎的生命。\par

% structure: p0102 source=h2 target=section reason=relative-heading-rank; base=h2

\section{第三十一章。}

从\Person{马克西米安}，天主——祂子民宗教的复仇者——将目光转向了那该诅咒的迫害的始作俑者\Person{加勒里乌}，为也藉对他的惩罚彰显祂威能的力量。\Person{加勒里乌}亦计划庆祝他的二十周年纪念；以此为由，他先前已用金银新税压榨各省，如今为以庆祝所许的节日来补偿他们，又用同样的借口重复压榨。谁能恰当地述说在征收赋税，尤其是谷物和其他地上果实方面，折磨人类所用的方法呢？所有不同官员的差役，或不如说刽子手，抓住每一个人，永不放脱。没有人知道该先向谁缴纳。对一无所有者也不予豁免，他们被要求在多种酷刑下立即缴纳，尽管无力支付。许多守卫四下安置，不给喘息之机，一年四季无一刻免于勒索。不同的官员，或不同官员的差役，常常为了向同一人收税的权利而争斗。没有打谷场没有收税人，没有葡萄园没有监视者，农夫所剩无几以维持生计！食物应从那些靠辛劳挣得它的人口中夺走，这已够痛苦；然而，有望事后得到缓解，或许能使这痛苦变得可忍受；但每个出现在纪念节日的人还须备办各式长袍，以及金银。有人或许会说：\Quote{无知暴君啊，你若夺走我土地的全部果实，强占预期的收成，我怎能置办这些东西？}就这样，在\Person{加勒里乌}的整个统治区域内，人们被剥夺财产，一切被搜刮进皇室国库，为的是皇帝能履行他庆祝一个他注定永远无法庆祝的节日的誓言。\par

% structure: p0104 source=h2 target=section reason=relative-heading-rank; base=h2

\section{第三十二章。}

\Person{马克西米努斯·达亚}因\Person{李锡尼}被擢升为皇帝而愤怒，他不愿再被称为凯撒，也不容许自己位居第三。\Person{伽勒里乌斯}屡次派人恳求\Person{达亚}让步，同意他的安排，尊重年长，敬重\Person{李锡尼}的白发。然而\Person{达亚}愈加傲慢。他声称，既然是他首先穿上紫袍，那么按占有权，他应有优先等级之位，并漠视\Person{伽勒里乌斯}的恳求和命令。那头畜生被刺痛，咆哮起来，因为那个他期望完全顺从而擢升为凯撒的卑鄙小人，竟忘了他所受的大恩，不虔地抗拒恩主的请求和意愿。\Person{伽勒里乌斯}终因\Person{达亚}的顽固而让步，废除了凯撒的副位称号，将奥古斯都称号授予自己和\Person{李锡尼}，将奥古斯都之子称号授予\Person{达亚}和\Person{君士坦丁}。\Person{达亚}稍后在致\Person{伽勒里乌斯}的信中趁机指出，在上次总集会上，他的军队曾以奥古斯都的称号向他欢呼。\Person{伽勒里乌斯}为此烦恼痛心，便命令四人皆应享有皇帝的称号。\par

% structure: p0106 source=h2 target=section reason=relative-heading-rank; base=h2

\section{第三十三章}

如今，当\Person{伽勒里乌斯}在位第十八年时，天主以不治之症击打了他。一处恶性溃疡在他的私处低处形成，并逐渐扩散。医生们试图根除它，将患处治愈。然而创口在愈合后再次破裂；一条血管爆裂，血液大量涌出，危及生命。尽管困难，血被止住了。医生们不得不重新手术，最终使伤口结疤。由于身体轻微活动，\Person{伽勒里乌斯}受到伤害，血流得比以前更多。他变得消瘦、苍白、虚弱，然后血被止住。溃疡对所用药物失去感觉，坏疽侵占了所有邻近部位。腐烂的肉被切除得越广，它扩散得越宽，用作治疗的一切手段反而加重了疾病。\par

\Quote{医术大师们退去了。}\par

于是从各处请来了名医；但任何人力手段都未成功。\Person{阿波罗}和\Person{阿斯克勒庇俄斯}被恳切祈求治愈良方：\Person{阿波罗}开了药方，疾病却加重了。已近致命危机，它占据了身体下部：他的肠子脱出，整个臀部腐烂。不幸的医生们虽然无望战胜疾病，仍不断地热敷和给药。体液被排斥后，疾病侵袭了他的肠子，体内生出了蛆虫。恶臭如此强烈，不仅充满宫殿，甚至弥漫全城；难怪那时他的膀胱和肠的通道已被蛆虫吞噬，变得不分，他的身体在难以忍受的痛苦中溶解成一团腐烂。\par

\Quote{痛彻灵魂，他因痛苦而咆哮，\newline{}受伤的公牛亦如此吼叫。}\Signature{Pitt}\par

他们将温热的动物肉敷在疾病的主要部位，希望温热能吸引出那些微小的蛆虫；因此，当敷料被移除时，涌出了无数蛆群：然而多产的疾病已孵化出更大量的蛆群来吞噬和消耗他的肠子。由于多种并发症，他的身体各部分已失去自然形状：上半部分干燥、瘦削、憔悴，他可怕的皮肤深陷于骨骼之间；而下半部分像膀胱一样肿胀，看不出关节的痕迹。这些事情在整整一年内发生；最终，被灾难所压倒，他不得不承认天主，在剧痛间歇中大声呼喊，他将重建他曾拆毁的教会，并为他的恶行赔补；当他临近死亡时，他发布了内容如下的谕令：——\par

% structure: p0112 source=h2 target=section reason=relative-heading-rank; base=h2

\section{第三十四章}

在我们为公共福祉的永久利益而制定的其他法规中，我们迄今致力于使一切符合罗马人的古老法律和公共纪律。\par

我们特别以之为目标，即那些抛弃了祖传宗教的基督徒，也应回归正确的见解。因为不知为何，这种刚愎和愚妄控制了它们，以致他们没有遵守那些可能是他们祖先确立的古老制度，反而任性为自己制定法律，并将许多信仰迥异的人聚集到不同的团体中。\par

在我们的谕令颁布后，命令基督徒使自己遵守古老制度，他们中许多人因恐惧危险而屈服，还有许多人陷入了危险；然而，因为仍有大量的人坚持自己的见解，并且我们发现目前他们既不对众神表示应有的敬畏和崇拜，也不崇拜他们自己的天主，因此，出于我们惯常的对所有人施予宽恕的仁慈，我们认为适合向那些人给予我们的宽恕，准许他们再次成为基督徒，并建立他们宗教集会的场所；但不得违反良好秩序。\par

我们拟通过另一道命令向官吏指明他们在此事上当如何行事。\par

\Quote{因此，基督徒因我们此番宽容，有责任向他们的天主祈祷，祈求我们的福祉、公共的福祉以及他们自身的福祉；使公共福祉在各方持续安全，他们自己也能安居于住所。}\par

% structure: p0118 source=h2 target=section reason=relative-heading-rank; base=h2

\section{第三十五章}

此诏令于五月朔日之前一日，在\Person{加莱里乌斯}第八次执政、\Person{马克西敏·达伊亚}第二次执政之时，颁布于\Place{尼科美底亚}。接着，监狱之门敞开，你，我至爱的\Person{多纳图斯}，偕同其他为信德作证的宣认者，从那座作了你们六年居所的牢狱中获得释放。然而，\Person{加莱里乌斯}并未因颁布此诏令而获得天主的宽恕。数日之后，他即被那引发全身腐烂的可怕疾病所吞噬。临终时，他将妻子和儿子托付给\Person{李锡尼}，交在了他手中。此消息在月底前已传至\Place{尼科美底亚}。他即位二十周年的庆典本应在随后的三月朔日举行。\par

% structure: p0120 source=h2 target=section reason=relative-heading-rank; base=h2

\section{第三十六章}

\Person{达伊亚}闻讯后，从东方以驿马疾驰，欲夺取\Person{加莱里乌斯}的领土，并趁\Person{李锡尼}滞留欧罗巴之际，擅自将整个地区直至\Place{迦克墩}海峡据为己有。进入\Place{比提尼亚}后，为迅速获取民心，他废除了\Person{加莱里乌斯}的税赋，众人皆大欢喜。两位皇帝之间出现纷争，几乎公开交战。他们各率军队对峙于两岸。然而，在特定条件下建立了和平与和睦。\Person{李锡尼}与\Person{达伊亚}在狭窄的海面上会晤，缔结条约，并握手以示友好。随后，\Person{达伊亚}以为万事安全，返回（\Place{尼科美底亚}），在其新领土上的所作所为与他在\Place{叙利亚}和\Place{埃及}时如出一辙。首先，他撤销了\Person{加莱里乌斯}给予基督徒的宽容和普遍保护，为此，他暗中指使各城呈递请愿书，要求在其城内不得建造任何基督教堂；他意图使自己的选择看上去像是被强求所致。应这些请愿书，他在宗教事务方面引入了一种新的治理模式，并为每座城市任命了一位大司祭，从最显赫的人士中选出。这些人的职责是每日向他们的所有神明献祭，并协助先前的司祭阻止基督徒建造教堂，或公开或私下崇拜天主；他授权他们强迫基督徒向偶像献祭，若遭拒绝，则将他们送交民事法官；而且，似乎这还不够，他在每个行省设立了一位总司祭，由国内首要人物之一担任；他命令所有这些新设立的司祭身穿白色礼服，这是衣着中最尊贵的标志。至于基督徒，他打算遵循他曾在东方采取的做法，并假作仁慈，禁止杀害天主的仆人，但下令对他们施加残害。于是，为信德作证的宣认者们被割裂耳朵和鼻孔，砍去手脚，挖出眼珠。\par

% structure: p0122 source=h2 target=section reason=relative-heading-rank; base=h2

\section{第三十七章}

正当他忙于这一计划时，他收到了\Person{君士坦丁}的书信，这封信阻止了他继续执行，因此他暂时隐藏了意图；然而，任何落入他手中的基督徒都被秘密投入海中。他也从未停止每日在宫殿中献祭的惯例。他还发明了一个做法：所有用作食物的动物，不是由厨师宰杀，而是由司祭在祭台前宰杀；因此，没有一样食物不经过品尝、祝圣和按异教礼仪洒酒就被端上；凡受邀赴宴的人，都必定污秽不洁地离开。在其他一切方面，他模仿他的师父\Person{加莱里乌斯}。因为如果有什么东西恰巧未被\Person{戴克里先}和\Person{马克西米安}动过，\Person{达伊亚}就会贪婪而不知羞耻地掠取。如今，每个人的谷仓都被关闭，所有仓库都被封存，尚未到期的税赋被提前征收。因此，由于耕作荒废，饥荒及一切物价飞涨。牲畜牧群被从牧场赶走用于每日献祭。他用祭肉的肉填饱士兵，使他们腐败，以至于他们不屑于通常的谷物配给，并随意将其丢弃。与此同时，\Person{达伊亚}赏赐他的众多卫兵以昂贵的衣袍和金质奖章，用银两犒赏普通士兵和新兵，并对在他军中服役的蛮族施以各种慷慨。至于他任意将活人的财产赐予宠臣，只要他们愿意索取属于他人的东西，我不知道这是否同样应得给予仁慈强盗的感谢，他们只抢夺而不杀人。\par

% structure: p0124 source=h2 target=section reason=relative-heading-rank; base=h2

\section{第三十八章}

但使他与众不同的，也是他超越所有前代皇帝之处，乃是他对淫乱妇女的欲望。我还能用什么词来形容它，除了说它是一种盲目而任性的情欲？然而，这些词语在我叙述他的暴行时，只能微弱地表达我的愤慨。罪行的重大使我的舌头无力，难以胜任其责。太监和淫媒四处搜寻，一旦发现任何美貌的面孔，丈夫和父母便被迫退避。贵妇和处女被剥去衣服，全身被检查，唯恐任何部分配不上皇帝的床榻。每当女子反抗，就被处以溺死；仿佛在这位通奸者的统治下，贞洁成了叛国。有些男人，看着自己因品德和忠诚而深爱的妻子被侵犯，无法忍受内心的痛苦，因而自杀。在这个怪物统治期间，唯有独特的丑陋才能保护任何女性的名誉免受他野蛮欲望的侵害。最后，他引入了一项习俗：禁止结婚除非得到皇帝许可；他以此作为满足其淫欲的工具。在奸污了自由出生的少女后，他将她们赐给他的奴隶为妻。他的同僚们也效仿皇帝的榜样，肆无忌惮地侵犯他们下属的妻室。因为谁能惩罚这些罪行呢？至于中等阶层人家的女儿，任何有意者都可以强行霸占。高贵的夫人，若不能强行霸占，则可请求皇帝作为赠礼获得。父亲也不能反对；因为一旦签署了皇帝的授权，他别无选择：要么死，要么接受某个蛮族人做他的女婿。因为禁卫军中几乎没有什么人，除了那些在\Person{戴克里先}第二十年被哥特人驱逐出家园，而屈服于\Person{加莱里乌斯}并进入他麾下的人。对人类而言，这是不幸的：那些从蛮族奴役中逃离的人，居然如此来统治罗马人。被这样的卫兵环绕，\Person{达亚}压迫并侮辱了东方帝国。\par

% structure: p0126 source=h2 target=section reason=relative-heading-rank; base=h2

\section{第三十九章。}

如今，\Person{达亚}在满足其淫欲时，以自己的意志作为是非的标准；因此他毫不克制地去追求\Person{加莱里乌斯}的遗孀、皇后\Person{瓦莱里娅}，他最近曾称她为母亲。她丈夫死后，她来到达亚处，因为她以为在他统治下生活比别处更安全，尤其因为他是个已婚男人；但这邪恶之徒立刻燃起了对她的欲望。\Person{瓦莱里娅}仍穿着丧服，服丧期尚未届满。他派人向她求婚，并承诺若她同意，他将休掉自己的妻子。她率直地作出了只有她敢作的回答：首先，她不会在穿着丧服且丈夫\Person{加莱里乌斯}的骨灰——他又是达亚的义父——尚未冷却时谈论婚事；其次，他提出休弃一位忠实的妻子以接纳另一位，而这位将来也可能被他抛弃，这是不敬之举；最后，对她这样身份和尊贵的女人而言，再次结婚是失礼、无先例且不法的。这大胆的回答传到\Person{达亚}耳中，他的欲望立刻化为愤怒和狂暴的怨恨。他宣布公主被剥夺财产，夺取她的物品，撤走她的侍从，将她的太监拷打至死，并将她和她的母亲\Person{普里斯卡}流放；但他并未指定她流放期间的特定居住地；因此，他侮辱性地将她在她漂泊途中任何落脚处都驱逐出去；而最甚者，他以虚假的通奸罪名判处那些与她最亲近、最信任的贵妇们死刑。\par

% structure: p0128 source=h2 target=section reason=relative-heading-rank; base=h2

\section{第四十章。}

有一位地位很高的贵妇，她已从多个儿子有了孙辈。\Person{瓦莱里娅}像对待第二位母亲那样爱她；\Person{达亚}怀疑她的建议导致了\Person{瓦莱里娅}拒绝他的求婚，于是他命令总督\Person{埃拉提努斯}以破坏她名誉的方式将她处死。除了她之外，还加上了另外两位同样高贵的妇人。其中一位，她的女儿在罗马是一名维斯塔贞女，她与流放的\Person{瓦莱里娅}保持着秘密联系。另一位，嫁给了一位元老，与皇后关系密切。卓越的美貌和品德成了她们死亡的原因。她们被拖到审判席前，那不是一个正直的法官，而是一个强盗。实际上也没有任何控告者，直到一个被指控其他罪行的犹太人，因希望得到赦免而被诱导，对无辜者作伪证。这位公平而警觉的行政长官派卫兵将他带出城外，以免群众用石头砸死他。这出悲剧在\Place{尼西亚}上演。犹太人被下令受刑，直到他说出被教唆的话，而行刑者用棍棒阻止妇女们为自己辩护。无辜者被判处死刑。于是，哀号哭泣四起，不仅陪同其贤妻的元老如此，在场的观众也是如此——他们被这丑闻未闻的审判吸引而来；为防止群众暴力从刽子手手中解救被定罪者，军事指挥官带着轻步兵和弓箭手跟随。就这样，在武装士兵的押送下，她们被带往刑场。她们的仆从被迫逃离，若不是朋友的怜悯秘密埋葬了她们，她们将得不到安葬。而对那个假扮的通奸者，赦免的承诺也没有兑现，因为他被绑上了绞刑架，随后他揭露了整个秘密阴谋；在最后一息，他向所有旁观者声明，这些妇女是无辜而死的。\par

% structure: p0130 source=h2 target=section reason=relative-heading-rank; base=h2

\section{第四十一章。}

然而，皇后被流放在\Place{叙利亚}某处荒僻之地，暗中将她所遭遇的灾祸告知其父\Person{戴克里先}。\Person{戴克里先}派遣使者到\Person{达亚}那里，请求将他的女儿送回。他未能成功。他一再恳求，但她仍未被送回。最后，他启用一位亲戚，此人是一位位高权重的军人，请他凭往日的恩情恳求\Person{达亚}。这位使者与其他人在交涉上同样无果，便向\Person{戴克里先}禀报，他的请求是徒劳的。\par

% structure: p0132 source=h2 target=section reason=relative-heading-rank; base=h2

\section{第四十二章。}

此时，根据\Person{君士坦丁}的命令，\Person{马克西米安·赫丘利}的雕像被推倒，其画像被移除；由于两位老皇帝的肖像通常画在同一幅画上，因而两幅肖像同时被移除。就这样，\Person{戴克里先}活着看到了任何先前的皇帝都未曾见过的耻辱，在精神烦恼与身体疾病的双重负担下，他决定死去。他辗转反侧，灵魂被忧伤搅动，既不能进食，也无法安歇。他时常叹息、呻吟、哭泣，不停地变换各种姿势，时而躺在卧榻上，时而又滚在地上。因此，这位曾二十年最为兴盛的皇帝，被贬为庶民的卑微境地，受到最粗暴的对待，被迫厌弃生命，最终无法接受营养，在精神的痛苦中耗尽而亡。\par

% structure: p0134 source=h2 target=section reason=relative-heading-rank; base=h2

\section{第四十三章。}

在天主的仇敌中，尚余一人，其倾覆与结局，我现在就要叙述。\par

\Person{达亚}从\Person{伽列里乌斯}优先选择\Person{李锡尼}时起，便对他心怀嫉妒与恶意；尽管他们最近缔结了和平条约，这些情绪依然存在。\Person{达亚}听说\Person{君士坦丁}的姐妹与\Person{李锡尼}订婚后，担心这两位皇帝通过缔结这一联姻，旨在联合对付他；于是他暗中派遣使者前往\Place{罗马}，请求与\Person{马克森提乌斯}结为友好联盟；他也以热忱的语气写信给他。使者受到有礼的接待，建立了友谊，作为标记，\Person{马克森提乌斯}与\Person{达亚}的肖像被公之于众。\Person{马克森提乌斯}欣然接受，仿佛这是来自天上的援助；因为他对\Person{君士坦丁}早已宣战，似乎是替其父\Person{马克西米安}之死复仇。从这种表面的孝道中产生了一种猜测：那个可憎的老头假装与儿子争吵，是为了有机会消灭他的权力对手，从而为自己和儿子铺平道路，以便拥有整个帝国。然而，这种猜测毫无根据；因为他的真正目的是要毁灭他的儿子和其他人，然后使他自己和\Person{戴克里先}复位至最高权力。\par

% structure: p0137 source=h2 target=section reason=relative-heading-rank; base=h2

\section{第四十四章。}

此时，\Person{君士坦丁}与\Person{马克森提乌斯}之间爆发了一场内战。尽管\Person{马克森提乌斯}留守\Place{罗马}，因为占卜者预言他若出城必将灭亡，但他仍通过能干的将领指挥军事行动。在兵力上他超过对手；因为他不仅有他父亲的军队（该军从\Person{塞维鲁}处投奔过来），还有他最近从\Place{毛里塔尼亚}和\Place{意大利}召集的自己的军队。他们交战，\Person{马克森提乌斯}的军队起初得胜。最后，\Person{君士坦丁}以坚定的勇气和准备应对一切事件的心志，率领全军前往\Place{罗马}附近，在\Place{米尔维安桥}对面扎营。\Person{马克森提乌斯}统治的周年纪念日临近，即十一月朔日前第六日，他统治的第五年即将结束。\par

\Person{君士坦丁}在梦中受指示，将天上的记号描绘在士兵的盾牌上，然后以此进行战斗。他照所吩咐的做了，在盾牌上刻下了字母\OriginalTerm{基督标记（凯乐符号）}{\GreekText{Χ}}，并在其中画一条垂直线，在顶部弯折，这就是基督的符号。有了这个记号（\GreekText{Χ}\par

\SourceRecord{Translated by William Fletcher.；Ante-Nicene Fathers；Vol. 7.；Edited by Alexander Roberts, James Donaldson, and A. Cleveland Coxe.；Buffalo, NY: Christian Literature Publishing Co.,；1886.；<http://www.newadvent.org/fathers/0705.htm>.}
